This thesis studies whole-cell instance segmentation of human hiPSC-derived cardiomyocytes in dense microscopy images with weak boundaries and strong cell-cell adhesion. We start from the released CellSAM checkpoint and perform an implementation-level audit to establish a reproducible segmentation branch and a unified final inference protocol. On top of this audited basis, we build a cardiomyocyte-specific adaptation framework with decoder-focused fine-tuning, biology-prior-aware channel handling, and explicit separation between Oracle and end-to-end evaluation. The sealed segmentation mainline is T28, reported by dual-seed mean on val71 (PQ@0.5 = 0.6837, BM-Dice = 0.8166), while locked test73 baselines are retained as a split-aware reference tier. For deployment, the selected detector-segmentation route is T33g(dapi\_cm, candidate-aligned, q35) + T28, framed as a biology-prior line rather than a detector-metric-optimal claim. The analysis indicates that the current deployment bottleneck is prompt quality rather than mask-decoder capacity. Overall, this work provides a domain-audited and prompt-aware adaptation framework for foundation segmentation in challenging cardiomyocyte imaging.
